% Full modern proofs for the exact new propositional presentations.
\begin{annexitem}{annex:p8.mi.operations}{thm:p08-mi-derivation-operations}{Derivation operations and exact base correspondence}
\begin{proof}
For an MPL tree, enlarging its permitted premise set leaves all actual
open leaves permitted. The tree is finite. Replace each selected open
leaf by a copy of its supplied derivation, first giving different copies
pairwise disjoint label sets fresh from the receiving tree. The old
inferences still have matching premises and discharge exactly their old
classes outside the inserted copies. If a class becomes empty, it is
legal on this sheet. The resulting open leaves are the unreplaced leaves
and the open leaves of the inserted copies; no other premise is introduced.

To obtain deduction, assign one fresh label to all open occurrences of
\(A\) to be discharged and apply implication introduction. If there are
none, use its allowed empty class. Conversely, apply implication
elimination to the derived \(A\to B\) and a new open \(A\)-leaf.
The conjunction and disjunction introduction/elimination operations are
the displayed rules. Reasoning by cases is disjunction elimination,
freshening the labels of the two minor trees before their indicated
classes are discharged. None of these arguments uses falsity elimination.

For any of the three Hilbert sequences, finite support and monotonicity
follow from their line definition. To replace a premise line, insert
the finite sequence deriving it from the new premises; later modus
ponens steps still refer to earlier available formulas. Only finitely
many premise lines occur. Thus replacement proves idempotence of
deductive closure as well as its extensivity and monotonicity.

Here is the Hilbert deduction construction, including its auxiliary
identity derivation. A1 gives \(A\to((A\to A)\to A)\) and
\(A\to(A\to A)\). The A2 instance with middle formula \(A\to A\)
and consequent \(A\), followed by two modus-ponens steps, gives
\(A\to A\). Now process a derivation from \(\Gamma\cup\{A\}\)
line by line. An axiom instance or \(\Gamma\)-premise \(D\) remains
available from \(\Gamma\); A1 and modus ponens give \(A\to D\).
For the premise \(A\), use the identity derivation. For modus ponens
from \(E\) and \(E\to D\), the earlier transformed lines are
\(A\to E\) and \(A\to(E\to D)\); A2 and two modus-ponens steps
give \(A\to D\). This covers also each calculus's extra axiom
instances. The reverse deduction direction is modus ponens with the
additional premise \(A\).

The conjunction schema with two modus-ponens steps gives conjunction
introduction, and the two projection schemas give its eliminations.
The injection schemas give disjunction introduction. For cases, use
deduction on the two minor derivations to get \(A\to C\) and
\(B\to C\), then apply the Cases schema and three modus-ponens steps
with \(A\lor B\). Falsity's schema and modus ponens give explosion.

We verify the stated base bridge rather than assuming a Hilbert/ND
identification. Each \(\mathsf H_0\) schema is a closed derivation in
the propositional empty-discharge NJ calculus. For A1, assume \(A\),
introduce \(B\to A\) with empty \(B\)-discharge, then discharge
\(A\). For A2, assume its antecedent \(A\to(B\to C)\), then
\(A\to B\), then \(A\); two applications give \(B\to C\) and
\(B\), followed by application to get \(C\). Discharge the three
assumptions in reverse order. The conjunction schema assumes \(A,B\),
forms their conjunction, and discharges \(B,A\); its projections assume
the conjunction, select the indicated constituent, and discharge it.
Each injection assumes its indicated disjunct, introduces the sum, and
discharges the assumption. For Cases, assume \(A\to C,B\to C,A\lor B\),
apply disjunction elimination using application in each minor branch,
then discharge the three outer assumptions in reverse order. Finally,
assuming \(\bot\), falsity elimination gives \(A\); discharge falsum.
Every binder and branch label in these trees is chosen fresh. Modus
ponens is NJ implication elimination, so induction over a Hilbert
sequence translates it from the same premise set into this NJ calculus.

Conversely, translate an empty-discharge propositional NJ tree by induction.
The conjunction, projection, injection, application and falsity cases use
the Hilbert operations already proved. In implication introduction, the
induction hypothesis supplies a Hilbert derivation from the surviving
premises together with \(A\); deduction internalizes \(A\), even if
its class is empty. For disjunction elimination, apply deduction to
each minor derivation and then the Hilbert Cases operation. All surviving
premises remain in the union of the original supporting sets.
This proves exact Hilbert/empty-discharge-NJ consequence equivalence.

To connect that NJ presentation to the specified propositional LJ relation,
one can use the following explicit consequence translations. An NJ
assumption translates to Id; conjunction introductions/projections,
injections, application and falsity elimination translate using the
sequent operations proved in
Lemma~\ref{thm:p08-ipl-finite-support-operations}; implication introduction
uses that lemma's deduction construction and disjunction elimination uses
its cases construction. Thus NJ gives LJ. For the reverse induction over
LJ, Id is an assumption; falsum left is NJ falsity elimination; WL and CL
change only allowed support. Conjunction left replaces its two constituent
assumptions by the two projections of a new conjunction assumption;
conjunction right uses NJ conjunction introduction. Disjunction left
uses cases and both right rules use injections. Implication left derives
its antecedent operand by the first minor derivation, applies the
implication assumption, then substitutes the consequent proof into the
second minor derivation. Implication right uses NJ introduction, allowing
empty discharge. Cut substitutes a proof into the distinguished
assumption occurrence; Mix substitutes separately freshened copies into
the indicated finitely many occurrences. These operations preserve
finite supporting sets. This covers every propositional LJ row, without
an equality or quantifier case. The imported lemma is used only for its
exact propositional LJ hypotheses; no minimal explosion is inferred from it.
Deleting falsity elimination from this NJ rule inventory is exactly the
new MPL sheet, proving the final assertion.
\end{proof}
\end{annexitem}

\begin{annexitem}{annex:p8.mi.minimal}{thm:p08-mi-minimal-adequacy}{Minimal falsum, soundness, prime extension and truth}
\begin{proof}
First every formula persists in a minimal model. Atoms persist by
assumption, and falsum persists because \(F_{\bot}\) is upward closed.
Conjunction and disjunction preserve persistence constituentwise. If
\(w\forces B\to C\) and \(v\geq w\), every \(u\geq v\) is also
above \(w\), so \(u\forces B\) implies \(u\forces C\).
This proves implication persistence and completes structural induction.

For soundness induct on an MPL tree, strengthening the assertion to
every world forcing all its actual open assumptions. An assumption is
forced by hypothesis. The conjunction introduction case uses both
induction conclusions; its two eliminations select the indicated
conjunct. Disjunction introduction forces the indicated disjunct at the
same world. Implication elimination uses reflexivity at that world.
In implication introduction, any future world forcing the discharged
formula also forces every surviving assumption by persistence. Apply
the induction hypothesis for the premise tree there; this gives the
consequent at every such future world. The argument works when the
discharged class is empty. In disjunction elimination, the induction
hypothesis for the major tree supplies one disjunct at the current world.
Apply the corresponding minor tree's induction hypothesis using that
disjunct and its surviving assumptions; either minor branch gives the
same conclusion. These are all the rules. There is deliberately no
falsity-elimination case and no assertion that a falsum world is impossible.

For completeness, let deductive closure mean closure under MPL
derivability. A minimal prime theory is a deductively closed set with
the disjunction property; no condition excluding \(\bot\) is imposed.
We now establish the needed extension construction. If a set avoids
\(C\), partially order its supersets avoiding \(C\) by inclusion.
The union of a chain still avoids \(C\), because a derivation uses
finitely many premises, all contained in one chain member. Zorn's lemma
therefore gives a maximal set \(T\) avoiding \(C\).
It is closed: if \(T\derives D\), adjoining \(D\) cannot derive
\(C\), because premise replacement would give \(T\derives C\).
Maximality thus puts \(D\) in \(T\). Also \(C\notin T\).
If \(B\lor D\in T\) but neither constituent belongs to \(T\),
maximality makes both \(T\cup\{B\}\) and \(T\cup\{D\}\) derive
\(C\). Reasoning by cases then makes \(T\) derive \(C\), a
contradiction. Thus \(T\) is prime. This argument did not use
explosion; \(\bot\) may belong to \(T\).

If a minimal prime theory \(T\) omits \(B\to D\), deduction implies
\(T\cup\{B\}\not\derives D\). Apply the extension construction
with target \(D\) to obtain a prime \(U\supseteq T\) containing
\(B\) but not \(D\). This is the failed-implication witness.

Given \(\Gamma\not\derives_{\mathsf{MPL}}A\), choose a prime
\(T_0\supseteq\Gamma\) omitting \(A\). The worlds are all minimal
prime theories extending \(T_0\), ordered by inclusion; the atom
valuation is membership, and
\[
 F_{\bot}=\{T:T_0\subseteq T,\ T\text{ is minimal prime},\ \bot\in T\}.
\]
This is a set-sized rooted partial order and all valuations, including
falsum, are persistent. The truth lemma is induction on formulas.
The atom and falsum cases are their membership definitions. Deductive
closure and conjunction introduction/projections give the conjunction
case in both directions. Disjunction introductions and primeness give
its two directions. If \(B\to D\in T\), each prime extension
containing \(B\) contains \(D\) by application and closure, giving
forcing of the implication. If it is omitted, the witness prime
extension just constructed forces \(B\) but not \(D\) by induction,
so the implication is not forced. Thus every formula is forced exactly
when it belongs to the world theory. At \(T_0\), all premises are
forced and \(A\) is not. This proves strong completeness, including
arbitrary premise sets. The proof is a classical ZFC argument;
set-sizedness follows because the worlds form a subset of the power
set of the formula set.
\end{proof}
\end{annexitem}

\begin{annexitem}{annex:p8.mi.falsum-recovery}{thm:p08-mi-minimal-falsum-recovery}{Falsum recovered without explosion}
\begin{proof}
From an open \(\bot\)-assumption, implication introduction with an
empty \(B\)-discharge gives \(B\to\bot\). On a separate copy of
the falsum assumption, implication introduction with empty \(\neg B\)-discharge gives
\(\neg B\to\bot\). Conjunction introduction gives
\(\neg B\land\neg\neg B\); the original falsum assumption
can now be discharged if an implication theorem is required.
Conversely, the two projections of that conjunction give
\(\neg B=B\to\bot\) and \(\neg\neg B=\neg B\to\bot\).
Applying the latter to the former gives \(\bot\). No step was
falsity elimination. Soundness transfers both directions to every
minimal model. If a world forces \(\bot\), it also forces
\(\bot\to\bot\), so it forces a formula and its negation.
If it forces \(D\) and \(D\to\bot\), reflexivity and the
implication clause give \(\bot\), proving the asserted world
characterization.

The translation replacing \(\bot\) by a new atom is a bijection
from the old formula grammar to the positive grammar over the enlarged
atom set. Its inverse replaces that distinguished atom by \(\bot\)
and leaves other atoms unchanged. It commutes with every binary constructor.
Applying it at every formula occurrence in an MPL tree preserves
all its rule instances, labels and selected discharge classes, since
none of those rules mentions a special falsum behavior. The inverse
map does the same. The translated positive ND rules have exactly
the positive Hilbert schemas in the base-bridge proof of
Annex~\ref{annex:p8.mi.operations}; remove just that proof's
falsity-schema/elimination case. Every remaining explicit schema
derivation and every ND induction case is unchanged, giving both
directions of the positive Hilbert equivalence. This establishes the
full grammar/derivability bridge without assuming that a primitive
negation symbol is an implication abbreviation.
\end{proof}
\end{annexitem}

\begin{annexitem}{annex:p8.mi.intermediate}{thm:p08-mi-intermediate-adequacy}{Greatest-world KC and linearly ordered LC completeness}
\begin{proof}
Formula persistence for these models is the proved IPL persistence lemma,
with the falsum set fixed empty. The base schemas are valid by the exact
base/NJ bridge and its rule soundness: the MPL soundness cases above
apply, and falsity elimination is valid because no world forces \(\bot\).
Modus ponens preserves truth at a common world. It remains to verify the
added schemas and the completeness frame restrictions.

On a greatest-world frame let \(t\) be greatest. For any formula \(B\),
if \(t\not\forces B\), persistence implies that no world forces
\(B\), so every world forces \(\neg B\). If \(t\forces B\), no
world can force \(\neg B\), since \(t\) is a future \(B\)-world
of each of them. Hence every world forces \(\neg\neg B\).
Classical case distinction gives weak excluded middle at every world.

For a linear frame, suppose neither \(B\to D\) nor \(D\to B\)
is forced at \(w\). The first failure supplies \(u\geq w\)
forcing \(B\) but not \(D\); the second supplies \(v\geq w\)
forcing \(D\) but not \(B\). If \(u\leq v\), persistence
contradicts the latter failure of \(B\); if \(v\leq u\), it
contradicts the former failure of \(D\). Linearity gives one of
these cases. Thus prelinearity holds. Induction over the respective
Hilbert sequences proves soundness from all their permitted premises.

For the following construction, a prime theory means a deductively closed,
consistent set with the disjunction property under the indicated Hilbert
relation; consistency means not deriving \(\bot\).
For either Hilbert extension, the finite-support, substitution, deduction
and cases operations in Annex~\ref{annex:p8.mi.operations} give the
same maximal-avoidance construction as for minimal logic. Every maximal
set avoiding \(C\) is closed and prime. This time it is also
consistent: if it derived falsum, explosion would derive \(C\).
Every prime theory contains all schema instances of its own calculus.
The failed-implication witness follows from deduction, as before.
The canonical prime cone above a root \(T_0\) therefore has the full
membership truth lemma, with falsum absent at every world. This verifies
the canonical model before imposing either special frame condition.

For LC, take any prime theories \(U,V\) above \(T_0\). If they
were incomparable, choose \(B\in U\setminus V\) and
\(D\in V\setminus U\). The prelinearity schema belongs to \(T_0\).
Primeness of \(T_0\) puts either \(B\to D\) or \(D\to B\)
there. The first belongs to \(U\), whence application gives
\(D\in U\), contradicting its choice. The second similarly
puts \(B\) in \(V\). Thus the entire canonical cone is linearly
ordered by inclusion. If LC fails to derive \(A\) from \(\Gamma\),
maximal avoidance supplies its root containing \(\Gamma\) and
omitting \(A\); the truth lemma gives the required linear countermodel.

For KC, first every consistent theory extends to a maximal consistent
theory by the chain-union/Zorn argument with target \(\bot\).
Such a maximal theory \(M\) decides every formula: if \(B\notin M\),
adjoining \(B\) derives \(\bot\), and deduction gives
\(\neg B\in M\). Both \(B\) and \(\neg B\) cannot belong
to a consistent theory, by application.

We show that the root prime theory \(T_0\) has exactly one maximal
consistent extension. For each formula \(B\), weak excluded middle
and primeness put \(\neg B\) or \(\neg\neg B\) in \(T_0\).
In the first case every consistent extension omits \(B\). In the
second every maximal consistent extension contains \(B\): otherwise
its decision property gives \(\neg B\), which with \(\neg\neg B\)
would derive falsum. Thus any two maximal consistent extensions of
\(T_0\) agree on every formula and are equal. Existence was already
proved; call the unique one \(M\).
Each prime \(U\supseteq T_0\) is consistent and has a maximal
consistent extension. That extension is also an extension of \(T_0\),
so it is \(M\), and \(U\subseteq M\). The maximal theory \(M\)
is itself prime by maximal avoidance, hence is a world in the canonical
cone and is its greatest world. Again a root avoiding the proposed
conclusion exists whenever consequence is not derivable, and the truth
lemma gives the required countermodel. No finiteness assumption was
used in either canonical-cone argument. All Zorn steps use classical
ZFC, as stated in the theorem.
\end{proof}
\end{annexitem}

\begin{annexitem}{annex:p8.mi.finite}{thm:p08-mi-finite-decision}{Filtration preserving falsum, a greatest world, and linearity}
\begin{proof}
Start with the respective rooted countermodel just proved. For the
finite subformula-closed set \(\Sigma\), put
\(\tau(w)=\{B\in\Sigma:w\forces B\}\). Identify worlds with
equal profiles and order the profiles by set inclusion. There are at
most \(2^{|\Sigma|}\) profiles. Inclusion makes their quotient a
partial order, and persistence makes the original root profile its root.
For each atom in \(\Sigma\), give a profile that atom exactly when
it contains it; atoms outside \(\Sigma\) receive empty valuations.
If \(\bot\in\Sigma\), let the filtered falsum upset be the
profiles containing it. If \(\bot\notin\Sigma\), choose that
upset empty. In KC and LC no original world forces falsum, so the
filtered falsum upset is empty in either case. The resulting valuations
are upward closed by profile inclusion.

Here is the truth induction, including the minimal falsum and implication
cases; IPL's false-at-every-world clause is not used for minimal falsum.
For an atom in \(\Sigma\), the valuation is exactly the original
truth bit. For falsum in \(\Sigma\), the new upset is exactly its
original truth bit. Conjunction and disjunction are preserved by the
induction hypotheses on their constituents, which are in \(\Sigma\).

Suppose \(B\to D\in\Sigma\) and \(w\forces B\to D\).
Every profile \(\tau(v)\supseteq\tau(w)\) contains that same
implication, so its chosen original representative forces it, even
when \(w,v\) were incomparable originally. If its filtered world
forces \(B\), induction gives \(v\forces B\); reflexivity and
the original implication give \(v\forces D\), and induction gives
filtered \(D\). This proves the filtered implication. Conversely,
if the original implication fails at \(w\), take an original future
\(v\) forcing \(B\) but not \(D\). Persistence gives
\(\tau(w)\subseteq\tau(v)\), and induction gives the same
witness in the filtered frame. Therefore the filtered implication
also fails. This proves truth preservation for every input subformula.

If the original KC frame has greatest world \(t\), persistence gives
\(\tau(w)\subseteq\tau(t)\) for every \(w\), so its filtered
profile is greatest. If the original LC frame is linear, any two
representatives are comparable; persistence makes their profiles
comparable, and the quotient order is linear. Thus the filter is a
countermodel in the exact required class, not merely an IPL filter.

The announced search is finite: for each size \(m\leq2^{|\Sigma|}\)
there are finitely many binary relations, roots, valuations of the finite
input atom set and falsum sets. Reflexivity, antisymmetry, transitivity,
rootedness, greatest-world/linearity conditions and upward closure are
decidable by finite inspection. Recursive evaluation of a finite formula
terminates, including the finite universal loop over accessible worlds
for implication. A found countermodel rules out derivability by
soundness. If none exists, the finite-countermodel result rules out
nonderivability and gives derivability. Thus both outcomes terminate.
There is no efficiency claim. This proof uses finite premises for the
finite \(\Sigma\); it does not filter an arbitrary infinite premise
set into a fixed finite bound.
\end{proof}
\end{annexitem}

\begin{annexitem}{annex:p8.mi.separations}{thm:p08-mi-separations}{Derivations, countermodels, and maximal versus greatest}
\begin{proof}
Adding falsity elimination gives the IPL consequence relation by the
base bridge. The one-world abnormal minimal model in the body refutes
\(\bot\to p\), which IPL derives. This proves the first strict
inclusion. Each intermediate sheet adds schemas to the same IPL base.

For \(\mathsf{KC}\subseteq\mathsf{LC}\), substitute
\(B=\neg A\) in prelinearity. The resulting two disjuncts imply
\(\neg A\) and \(\neg\neg A\), respectively, in the common
intuitionistic base. Indeed, under \(A\to\neg A\), assume \(A\),
apply the premise to obtain \(\neg A\), apply it to \(A\) to
obtain falsum, and discharge \(A\) to get \(\neg A\).
Under \(\neg A\to A\), assume \(\neg A\), apply the premise
to obtain \(A\), and apply the assumed negation to obtain falsum;
discharge that assumption to get \(\neg\neg A\). These are NJ
derivations and hence Hilbert-base derivations by the proved bridge.
Cases and injections now derive weak excluded middle in LC. Replacing
each KC additional-axiom line by this LC derivation translates any KC
theorem to LC. All shared base schemas and prelinearity are classically
valid, and modus ponens is classically sound; thus LC is contained in CPC.

We give the three remaining strictness witnesses. A root \(r\) with
two incomparable terminal successors \(u,v\), with \(p\) true just
at \(u\), is an IPL model. At \(r\), \(\neg p\) fails because
\(u\forces p\); \(\neg\neg p\) fails because
\(v\forces\neg p\). Thus KC's weak excluded middle is not an IPL
theorem. Next take \(r<a<t\) and \(r<b<t\), with \(a,b\)
incomparable, put \(V(p)=\{a,t\}\) and \(V(q)=\{b,t\}\).
This frame has greatest world \(t\), so is a KC model. At \(r\),
\(p\to q\) fails at \(a\) and \(q\to p\) fails at \(b\);
LC prelinearity is therefore not a KC theorem. Finally a two-world
chain \(r<t\) with \(V(p)=\{t\}\) is an LC model. The root
forces neither \(p\) nor \(\neg p\), so excluded middle, a CPC
theorem, is not an LC theorem. Soundness makes each witness a
nonderivability proof for the stated calculus.

For KC's disjunction property, its schema \(\neg p\lor\neg\neg p\)
is a theorem, but neither disjunct is. The one-world Boolean model
with \(p\) true refutes \(\neg p\); the one-world Boolean model
with \(p\) false refutes \(\neg\neg p\). Both are greatest-world
models. For LC, prelinearity's instance
\((p\to q)\lor(q\to p)\) is a theorem. The Boolean valuation
\(p=1,q=0\) refutes the first disjunct; the reversed valuation
refutes the second. These one-world models are linear. The two failures
are theorem-level disjunction-property failures and do not contradict
the separate IPL disjunction theorem.

In a finite rooted partial order every world lies below a maximal world:
if it is not maximal, move strictly upward, and a strictly increasing
sequence cannot continue longer than the finite world set. If the
maximal world is unique, every world is below it, so it is greatest.
A greatest world is necessarily the unique maximal world. If there
are two distinct maximal worlds \(u,v\), set \(V(p)=\{u\}\),
an upset because \(u\) is maximal. At the root, \(\neg p\)
fails at \(u\) and \(\neg\neg p\) fails at \(v\), which forces
\(\neg p\). Thus universal validity of weak excluded middle forces
uniqueness of the maximal world and hence greatestness. The converse
is the schema soundness proof in Annex~\ref{annex:p8.mi.intermediate}.

For the infinite boundary take a root \(r\), a terminal world \(a\)
above it, and a separate chain \(r<b_0<b_1<\cdots\), with \(a\)
incomparable with every \(b_n\). The only maximal world is \(a\),
but it is not greatest. Put \(V(p)=\{a\}\). The root fails
\(\neg p\) because of \(a\) and fails \(\neg\neg p\)
because every \(b_n\) forces \(\neg p\). This supplies the
announced explicit failure of the finite equivalence when finiteness
is dropped.
\end{proof}
\end{annexitem}

\begin{annexitem}{annex:p8.mi.kc-algebras}{thm:p08-mi-kc-algebra-bridge}{The new-bottom construction and its finite frame representation}
\begin{proof}
Adding a new bottom to a distributive lattice keeps distributivity:
identities using only old elements are unchanged, and an occurrence
of the new bottom reduces using its meet/join equations. In any finite
distributive lattice, the elements \(z\) satisfying \(x\wedge z\leq y\)
are closed under joins. Their finite join is their greatest element,
so the relative pseudocomplement exists. Thus \(\alpha^*\) is a
Heyting algebra. For old \(x,y\), its implication agrees with that
in \(\alpha\): the old greatest admissible element is still
admissible and still above every old candidate and the new bottom.
The meet of two old elements is old, hence strictly above \(f\).
Consequently an old \(x\) has \(x\to f=f\), while \(f\to f=1\).
Weak excluded middle therefore takes top in \(\alpha^*\), but we
still need the exact semantic bridge for arbitrary formulas.

We supply the finite representation needed for that bridge. In a finite
Heyting algebra \(H\), a proper lattice filter contains top, is upward
closed and closed under finite meets, and omits bottom. It is prime
when membership of a join entails membership of a constituent.
If a filter \(P\) does not contain an element \(b\), and the
filter generated by \(P\cup\{a\}\) does not contain \(b\),
choose a maximal filter containing those data but omitting \(b\);
finiteness ensures its existence. It is prime. Indeed, if neither
\(x\) nor \(y\) belongs to it but \(x\vee y\) does, maximality
gives \(p\wedge x\leq b\) and \(q\wedge y\leq b\) for
some \(p,q\) in the filter. Distributivity then gives
\(p\wedge q\wedge(x\vee y)\leq b\); upward closure would put
\(b\) in it, a contradiction. This argument also applies to the
initial filter generated by any element \(a\not\leq b\), which
omits \(b\); hence prime filters separate \(a\not\leq b\).

Order all prime filters by inclusion and send each element \(a\)
to the upset of filters containing it. Separation makes this map
injective; membership gives preservation of bottom/top and meet,
and primeness gives join. For implication, if \(a\to b\in P\),
every prime extension containing \(a\) contains \(b\), because
\(a\wedge(a\to b)\leq b\). If \(a\to b\notin P\), the
filter generated by \(P\cup\{a\}\) omits \(b\): otherwise
\(p\wedge a\leq b\) for some \(p\in P\), so
\(p\leq a\to b\), putting that implication in \(P\).
The preceding maximal-filter argument gives a prime extension
containing \(a\) and omitting \(b\). This proves exact preservation
of the Kripke implication clause. Induction therefore preserves
every formula evaluation under any algebra valuation.

For \(H=\alpha^*\), its nonbottom elements form a prime filter:
their meet is old and nonbottom, and a nonbottom join has a nonbottom
constituent. This filter contains every proper prime filter and
therefore is the greatest world of the finite prime-filter frame.
At each evaluation world its future cone is rooted and still has
that greatest world. KC soundness thus applies at every prime filter.
If every premise has top algebra value, it holds in every prime
filter; soundness puts the conclusion in every prime filter.
Separation from top proves its value is top. This establishes the
sound algebraic direction, including the degenerate one-element
\(\alpha\) case after its nontrivial two-element extension.

For the converse, take a finite rooted greatest-world KC countermodel
\(P\) supplied by the finite-countermodel theorem. Its upset lattice
has bottom \(\varnothing\), top \(P\), and the Heyting implication
\(U\Rightarrow V=\{w:\uparrow w\cap U\subseteq V\}\).
That implication is an upset and is the greatest upset \(Z\) with
\(U\cap Z\subseteq V\): one direction follows by testing
\(w\in U\cap Z\), and the other by persistence of \(Z\)
at every future \(U\)-world. This proves the lattice is Heyting
without an imported representation assumption. Since \(P\) has a
greatest world \(t\), every nonempty upset contains \(t\).
The nonempty upsets form a finite distributive lattice \(\alpha\)
with bottom \(\{t\}\) and top \(P\); it has a Heyting conditional
by the finite-distributive argument above. Adjoining the empty upset
as a new bottom reconstructs the entire upset algebra as \(\alpha^*\).
Assign each atom its truth upset. The operations coincide with the
forcing clauses, so induction gives every formula its truth set.
Each premise is forced at the root and therefore has top value;
the conclusion is not forced there and has a nontop value.
This is an algebraic countervaluation of the required form.
The input was finite exactly where the finite-countermodel theorem
was used; no arbitrary-premise finite-algebra completeness is inferred.
\end{proof}
\end{annexitem}

\begin{annexitem}{annex:p8.mi.chains}{thm:p08-mi-finite-chain-bridge}{Finite upsets, Godel operations, and the finite-input restriction}
\begin{proof}
Let \(W\) be a finite rooted linear frame. Its upsets form a finite
bounded chain under inclusion, with bottom \(\varnothing\), top
\(W\), meet intersection and join union. A formula's truth set is an
upset by persistence. For upsets \(U,V\), its implication truth set
is \(\{w:\uparrow w\cap U\subseteq V\}\).
If \(U\subseteq V\), this is \(W\). If \(U\not\subseteq V\),
linearity gives \(V\subsetneq U\). Every \(w\in V\) satisfies
the condition, by upward closure of \(V\). If \(w\notin V\),
choose \(u\in U\setminus V\). If \(w\leq u\), that \(u\)
is an accessible witness to failure. If \(u\leq w\), upward
closure of \(U\) puts \(w\) in \(U\setminus V\), and \(w\)
itself is a witness. Thus the implication truth set is exactly \(V\).
These are precisely the displayed Godel-chain operations. The map
from formulas to their truth sets therefore respects every constructor,
by structural induction. Root truth is exactly top value, because
an upset contains the least world iff it is the whole frame.

Conversely take any nontrivial finite bounded chain with elements
\(c_0<\cdots<c_n\), where \(c_0=0,c_n=1\). A linearly ordered
frame with \(n\) worlds has exactly \(n+1\) upsets. Map the chain
in order to those upsets, including the empty and full upsets.
This bijection preserves meet, join, bottom, top and the conditional
operation just computed. Assign each atom its corresponding upset.
Structural induction reproduces all chain evaluations as truth sets.
The one-element algebra supplies no counterexample because every
formula has top value, so excluding it here loses no failure case.
This proves equivalence between finite-chain countervaluations and
finite-linear-frame countermodels. The finite countermodel theorem and
LC soundness give the first assertion for every finite input.

Now fix an infinite bounded linearly ordered set. Any finite chain embeds
order-preservingly into it with endpoints fixed: choose finitely many
distinct interior elements and order them. Such a choice exists for
each required finite size because the fixed chain is infinite; density
is not required. The embedding preserves minimum, maximum and the
displayed conditional, whose value is either top or its second argument.
Thus it takes each finite-chain countervaluation to one in the fixed
infinite chain without changing which input values are top.

In the reverse direction, a valuation of a finite input uses only
finitely many atom values. Together with \(0,1\) they form a finite
subchain closed under all three operations: min/max select an operand,
and conditional returns top or its second operand. Induction shows
that all input formulas take exactly their original values in this
finite subchain. A countervaluation in the infinite chain therefore
already is one in a finite chain. This proves the second finite-input
equivalence, and proves precisely why no arbitrary-premise
fixed-numerical-chain claim follows: an infinite premise set need not
use only finitely many atom values, so this finite-subchain step is
not available. The result does not specify Dummett's denumerable
carrier or reconstruct his original argument.
\end{proof}
\end{annexitem}
